%% chapter1.tex: Introduction

\chapter{Introduction}
\label{chap:intro}

This chapter introduces the context and motivation for this work in
Section~\ref{sec:context}, states the problem addressed in
Section~\ref{sec:problem}, defines the objectives in
Section~\ref{sec:objectives}, lists the contributions in
Section~\ref{sec:contributions}, and outlines the document structure in
Section~\ref{sec:structure}.

%% -----------------------------------------------------------------------
\section{Context and Motivation}
\label{sec:context}

The deployment of autonomous mobile robots in environments where human presence
is dangerous, impractical, or impossible has grown considerably in recent years.
Search and rescue missions in disaster zones, industrial inspections in
hazardous facilities, and surveillance operations in areas without fixed
infrastructure all share a fundamental requirement: the robot must maintain a
reliable communication link with a remote operator, even when the two are not
within direct radio range. The absence of cellular towers, WiFi access points,
or any fixed networking equipment means that the communication network must be
created and maintained dynamically by the nodes themselves.

In such scenarios, Mobile Ad-hoc Networks (MANETs) emerge as the enabling
technology~\cite{conti2014mobile}. Each node becomes simultaneously a data
source, a data destination, and a potential relay that forwards packets on
behalf of others. This multi-hop communication paradigm extends the effective
range far beyond what direct wireless links could achieve: a robot operating
deep inside a building can stream video to an operator outside by having one
or more intermediate nodes relay each packet hop-by-hop until it reaches its
destination, even when source and destination lack a direct radio path.

A defining characteristic of robotic MANETs is the rapid and continuous
topology change caused by node mobility. Unlike sensor networks where devices
remain stationary after deployment, mobile robots navigate freely through their
environment following mission requirements. A link that exists at one moment
may vanish seconds later when a robot turns a corner, enters a different room,
or simply moves beyond radio range. Routing (the process of determining
paths for data to flow from source to destination) cannot rely on static
precomputation but must adapt continuously as the topology evolves.

A further challenge arises when the application traffic is a real-time video
stream. Video streaming imposes strict constraints on latency, jitter, and
packet loss: degradation below a perceptual threshold renders the stream
unusable for teleoperation, with direct consequences for the safety and
effectiveness of the mission. Routing adaptation must therefore occur quickly
enough to avoid prolonged stream interruptions when a link fails and an
alternative relay path must be activated.

Time Division Multiple Access (TDMA) provides an effective foundation for
organising transmissions in the shared wireless medium. By assigning each node
an exclusive time slot within a periodic frame, collisions are avoided by
construction and channel access is predictable. The RA-TDMAs+ protocol
developed at FEUP~\cite{oliveira2018clockless} extends this model to ad-hoc
mesh networks with relative synchronisation (that is, without requiring a
global clock) and distributes topology information via periodic beacon
packets. An experimental implementation of RA-TDMAs+ was realised by Diogo
Almeida~\cite{almeida2023synchronization}, providing a functional TDMA
framework as the foundation upon which the present work is built.

%% -----------------------------------------------------------------------
\section{Problem Statement}
\label{sec:problem}

The RA-TDMAs+ framework provides collision-free medium access and distributes
a connectivity matrix among nodes via periodic beacons. What it does not
provide is a routing layer: a mechanism that translates topology information
into forwarding decisions that allow arbitrary IP traffic to traverse multi-hop
paths transparently.

Several specific problems arise when adding such a routing layer to a TDMA
mesh network for robotic applications. First, the routing mechanism must
operate within the strict timing constraints imposed by the TDMA slot
structure, computing and installing routes without disrupting the medium-access
schedule. Second, it must detect link failures rapidly and trigger relay
activation before the ongoing application stream (in this case a video feed)
is perceived as interrupted by the operator. Third, the forwarding of
transit packets must be implemented efficiently: user-space forwarding planes
introduce additional latency and complexity, while kernel-level approaches
offer proven performance, standard observability, and seamless integration with
the operating system's networking stack. Fourth, a metric is needed to quantify link quality
within the TDMA context, since classical metrics such as hop count or received
signal strength do not capture the slot-level packet delivery behaviour that is
most relevant here.

This dissertation addresses these problems through the design and experimental
validation of a routing mechanism for RA-TDMAs+ networks, evaluated across
multiple network configurations and topology conditions to assess routing
robustness and adaptability.

%% -----------------------------------------------------------------------
\section{Objectives}
\label{sec:objectives}

The primary objective of this dissertation is to design, implement, and
evaluate a routing mechanism for a MANET built on top of the RA-TDMAs+
framework, enabling peer-to-peer IP communication between any two nodes in
the network independently of the current topology.

This primary objective decomposes into the following specific objectives:

\begin{enumerate}

    \item Define a \textbf{link quality metric} that quantifies the reliability
    of each wireless link based on TDMA beacon reception statistics, and use it
    to drive both Minimum Spanning Tree (MST) computation and relay triggering;

    \item Implement a \textbf{topology-aware routing algorithm} that derives
    next-hop decisions from an MST computed over the link quality graph
    disseminated by the MATRIX beacon protocol, and installs the resulting
    routes into the Linux kernel routing table;

    \item Integrate the relay forwarding mechanism with the \textbf{Linux
    kernel networking stack} using a dual routing table architecture, so that
    transit packets are forwarded by the kernel's own IP forwarding engine
    without user-space intervention;

    \item Implement a \textbf{video quality feedback loop} in which the base
    station signals video loss to the robot node, triggering immediate link
    quality degradation and route recomputation to activate the relay;

    \item \textbf{Validate} the complete system through practical
    experimentation, measuring routing convergence time after a link failure,
    end-to-end round-trip latency, video stream inter-arrival regularity, and
    throughput under both direct and relay communication conditions.

\end{enumerate}

%% -----------------------------------------------------------------------
\section{Contributions}
\label{sec:contributions}

The work reported in this dissertation makes the following contributions:

\begin{itemize}

    \item \textbf{Link quality metric for RA-TDMAs+}: a score in the range
    0--100 computed per TDMA slot from beacon reception statistics, integrating
    packet delivery rate estimation and consecutive slot miss detection, used
    to weight MST computation and to trigger relay activation;

    \item \textbf{MST-based routing layer}: design and implementation of a
    routing algorithm that builds a Minimum Spanning Tree over the link quality
    graph provided by the MATRIX protocol, derives a per-destination next-hop
    table from the spanning tree structure, and installs the computed routes
    into the Linux kernel routing table via Netlink;

    \item \textbf{Dual routing table architecture for relay forwarding}: a
    Linux kernel integration that uses two routing tables: one for traffic
    originating at the local node, transported through a TUN virtual interface
    over the TDMA channel, and one for transit traffic identified by source IP
    address and forwarded directly via the physical Wi-Fi interface using the
    kernel's built-in IP forwarding, eliminating the need for a user-space
    relay plane;

    \item \textbf{Video quality feedback loop}: a closed-loop mechanism in
    which the base station detects video stream loss and sends a UDP feedback
    signal to the robot node, which immediately sets the corresponding link
    quality to zero and schedules route recomputation, triggering relay
    activation without waiting for TDMA topology convergence;

    \item \textbf{Experimental evaluation}: a test methodology and measurement
    scripts for characterising routing convergence time, round-trip latency,
    video stream inter-arrival regularity, and throughput in a three-node
    scenario with topology changes induced by controlled signal attenuation.

\end{itemize}

%% -----------------------------------------------------------------------
\section{Document Structure}
\label{sec:structure}

The remainder of this document is organised as follows:

\begin{itemize}

    \item \textbf{Chapter~\ref{chap:background}} presents the theoretical
    background relevant to this work. It covers the characteristics and
    challenges of MANETs, introduces the RA-TDMAs+ TDMA framework and its
    beacon-based topology dissemination mechanism, reviews routing algorithms
    for ad-hoc networks with emphasis on MST-based approaches, and analyses
    related work on Layer 2 routing in RA-TDMAs+ networks.

    \item \textbf{Chapter~\ref{chap:implementation}} describes the design and
    implementation of the proposed system. It covers the MATRIX protocol and
    the link quality metric, the MST computation and next-hop derivation, the
    Netlink-based kernel route installation, the dual routing table architecture
    for relay forwarding, the video quality feedback mechanism, and the thread
    model that coordinates these components within the TDMA slot structure.

    \item \textbf{Chapter~\ref{chap:evaluation}} presents the experimental
    evaluation. It describes the hardware platform, the three-node test
    scenario and the methodology used to induce and measure topology changes,
    and discusses the results obtained for convergence time, round-trip
    latency, video stream inter-arrival regularity, and throughput under direct
    and relay communication.

    \item \textbf{Chapter~\ref{chap:conclusion}} summarises the work, draws
    conclusions from the experimental results, and identifies directions for
    future research.

\end{itemize}
